\documentclass[12pt]{article}
\usepackage[margin=0.3in,top=0.25in,bottom=0.4in]{geometry}
\makeatletter
\renewcommand{\@maketitle}{\newpage\null\vskip -3.5em\begin{center}{\LARGE \@title \par}\end{center}\par\vskip 0.2em}
\makeatother
\usepackage{amsmath,amssymb}
\usepackage{booktabs}
\usepackage{graphicx}
\usepackage{hyperref}
\setlength{\parskip}{0.6em}
\setlength{\parindent}{0pt}
\newcommand{\pt}{p_\perp}

\title{The partonic exclusive cross section at large $p_\perp$: GBW and MV}
\author{}
\date{}

\begin{document}
\maketitle

\section*{In short}

Everything here is the \textbf{partonic} cross section $\gamma A\to c\bar c A$, differential in the
transverse momentum $p_\perp$ of the charm quark, at fixed $z=p^+/q^+$ and fixed $x_{\mathbb P}$: no photon flux
and no fragmentation. In the code this is the function\\
\texttt{exclusiveCrossSection\_fixed\_qp}.

\begin{itemize}
\item The exclusive process $\gamma A\to c\bar c A$ has one $q\bar q$ dipole in the amplitude, so it
uses the dipole amplitude in the \textbf{fundamental} representation, $\mathcal{N}_F$. In the code this
is \texttt{dipole->N} (\texttt{src/integrands.cpp}); the $q\bar qg$ (TMD) part uses the adjoint one,
\texttt{dipole->N\_A} $=2\mathcal{N}-\mathcal{N}^2$.
\item At large $\pt$ the integrals are controlled by small dipoles, $r\sim1/\pt$. Only the parts of the
integrand that are \textbf{not smooth} at $r=0$ (logarithms) contribute to the large-$\pt$ tail: the
smooth parts give zero. This is the colour-transparency cancellation in a technical form.
\item For a dipole $\mathcal{N}=Q_{s0}^2r^2/4$ (GBW at small $r$) the only logarithm is the one of
$K_\nu(mr)$, and
\[
  I_0\to-\frac{Q_{s0}^2}{\pt^4},\qquad I_1\to\frac{2mQ_{s0}^2}{\pt^5}
  \qquad\Longrightarrow\qquad \frac{d\sigma}{d^2p\,dy}\propto\frac{m^2Q_{s0}^4}{\pt^8}.
\]
\item For the MV dipole, $\mathcal{N}=\tfrac{Q_{s0}^2r^2}{4}\ln\tfrac{1}{r\Lambda}$, the extra logarithm
gives $I_1\to Q_{s0}^2/(2m\pt^3)$, without the mass suppression, and
\[
  \frac{d\sigma}{d^2p\,dy}\propto\frac{(z^2+\bar z^2)\,Q_{s0}^4}{4\,\pt^6}
  +\frac{m^2Q_{s0}^4\,\ln^2(\pt^2/m\Lambda)}{\pt^8}.
\]
So the true asymptotic power for MV is $\pt^{-6}$, reached only above about $20$~GeV.
\item Our dipole (MV$^e$, $e_c=18.9$) has a perturbative tail at all $k\gtrsim3$~GeV: $k^{-3}$ below
$1/r_0=e_c\,e\,\Lambda\approx12$~GeV (from a non-smooth $r_0|\mathbf r|$ term) and the MV $k^{-4}$ above.
In the cross section this tail enters with a weight $\propto(r_0\pt)^2$, so in our range the $m^2/\pt^8$
term still dominates: local power $-7.3$ to $-7.5$ for $5<\pt<20$~GeV.
\item rcBK evolution builds a strong $r^2\ln\frac1r$ term (three times the MV one at $x_{\mathbb P}=10^{-3}$),
so at fixed small $x_{\mathbb P}$ the spectrum is close to $\pt^{-6}$. In the real
kinematics $x_{\mathbb P}\propto m_T$ grows with $\pt$, the evolution fades, and at $y=0$ the partonic
power is $-6.7$ at 5~GeV and $-8.2$ at 12~GeV (Section 7c).
\end{itemize}

\section*{1. The cross section and the fundamental dipole}

\textbf{Amplitude.} The photon splits into a $c\bar c$ pair of transverse size $\mathbf{r}$, the pair
scatters elastically on the target, and we project on a quark with transverse momentum $\mathbf{p}$:
\[
  \mathcal{A}_{h\bar h}(\mathbf{p})\;\propto\;\int d^2\mathbf{r}\;e^{-i\mathbf{p}\cdot\mathbf{r}}\;
  \Psi^{\gamma\to q\bar q}_{h\bar h}(\mathbf{r})\;\mathcal{N}_F(\mathbf{r}).
\]
The pair is a colour triplet--antitriplet, so it sees the target through the fundamental Wilson lines
$V$:
\[
  \mathcal{N}_F(\mathbf{r})=1-\frac{1}{N_c}\big\langle\mathrm{tr}\,V(\mathbf{x})V^\dagger(\mathbf{y})\big\rangle,
  \qquad \mathbf{r}=\mathbf{x}-\mathbf{y}.
\]
(In the $q\bar qg$ contribution the soft gluon is the one that scatters, so the adjoint dipole appears;
in the Gaussian/MV approximation at large $N_c$, $1-\mathcal{N}_A=(1-\mathcal{N}_F)^{C_A/C_F}\simeq(1-\mathcal{N}_F)^2$,
i.e.\ $\mathcal{N}_A=2\mathcal{N}_F-\mathcal{N}_F^2$, as in the code.)

\textbf{Photon wavefunction at $Q^2=0$.} There are two helicity structures:
\begin{itemize}
\item opposite quark helicities: $\Psi\propto(z\ \text{or}\ \bar z)\;\boldsymbol{\epsilon}\cdot\hat{\mathbf r}\;
m\,K_1(mr)$,
\item equal quark helicities (needs the mass): $\Psi\propto m\,K_0(mr)$.
\end{itemize}

\textbf{Angular integral.} With $\mathcal{N}_F$ depending only on $r=|\mathbf r|$,
\[
  \int_0^{2\pi}\!d\phi\;e^{-ipr\cos\phi}=2\pi J_0(pr),\qquad
  \int_0^{2\pi}\!d\phi\;\hat{\mathbf r}\,e^{-ipr\cos\phi}=-2\pi i\,\hat{\mathbf p}\,J_1(pr).
\]
So both amplitudes are proportional to $m$ times
\begin{equation}
  I_\nu(p)=\int_0^\infty dr\;r\,J_\nu(pr)\,K_\nu(mr)\,\mathcal{N}_F(r),\qquad\nu=0,1,
  \label{eq:I}
\end{equation}
and squaring and summing the helicities gives the formula of \texttt{integrand\_exclusive\_*}:
\begin{equation}
  \frac{d\sigma}{d^2\mathbf p\,dy}\;\propto\;z\,m^2\Big[(z^2+\bar z^2)\,I_1^2+I_0^2\Big].
  \label{eq:sigma}
\end{equation}
Everything below is about the large-$p$ behaviour of $I_0$ and $I_1$.

\section*{2. The tool: large $p$ from small $r$}

$K_\nu(mr)$ decays exponentially at large $r$, so the large-$p$ tail of Eq.~\eqref{eq:I} is fixed by
the behaviour of the integrand at small $r$. Expand the integrand in powers $r^a\ln^kr$ and use
\begin{equation}
  M_\nu(a)\equiv\int_0^\infty dr\;r^{a}J_\nu(pr)
  =\frac{2^{a}}{p^{a+1}}\,\frac{\Gamma\!\big(\tfrac{1+\nu+a}{2}\big)}{\Gamma\!\big(\tfrac{1+\nu-a}{2}\big)},
  \qquad
  \int_0^\infty dr\;r^{a}\ln^k\!r\;J_\nu(pr)=\partial_a^k M_\nu(a).
  \label{eq:mellin}
\end{equation}
(The integral is defined by analytic continuation in $a$; this is the standard Mellin-transform
method for asymptotic expansions.)

\textbf{Key fact: smooth terms give zero.} For $\nu=0$ and $a=1,3,5,\dots$ the denominator is
$\Gamma(0),\Gamma(-1),\Gamma(-2),\dots$, so $M_0(a)=0$. For $\nu=1$ the same happens for $a=2,4,6,\dots$
These are exactly the terms that are smooth functions of the 2D vector $\mathbf r$ (with the angular
factor $e^{i\nu\phi}$), and the Fourier transform of a smooth function falls faster than any power.
\emph{Only non-smooth terms, here the logarithms, produce a power-law tail.}

\textbf{Why.} $I_\nu$ is a 2D Fourier transform,
$I_\nu(p)\propto\int d^2\mathbf r\,e^{-i\mathbf p\cdot\mathbf r}f(\mathbf r)$ with
$f(\mathbf r)=K_\nu(mr)\,\mathcal N(r)\,e^{i\nu\phi}$.
\begin{itemize}
\item If $f$ is $C^\infty$ and decays fast, every integration by parts gives one more $1/p$, so the
transform falls faster than any power. A $p^{-n}$ tail therefore needs a point where $f$ is not smooth.
At large $r$, $K_\nu\sim e^{-mr}$ removes everything, and the only possible non-smooth point is $r=0$.
This is the precise meaning of ``large $p$ is controlled by $r\sim1/p$''.
\item ``Smooth'' means smooth in the vector $\mathbf r=(x,y)$, not in $r$. For $\nu=0$,
$r^2=x^2+y^2$, $r^4,\dots$ are polynomials. For $\nu=1$, $re^{i\phi}=x+iy$, $r^3e^{i\phi},\dots$ are
polynomials too. Their Fourier transforms are derivatives of $\delta^{(2)}(\mathbf p)$, so they vanish
at $p\neq0$. With the measure $r\,dr$ these are exactly the zeros of $M_\nu(a)$ above. Instead
$r^2\ln r$ (whose second derivatives diverge as $\ln r$) and $r^{2\gamma}$ with non-integer $\gamma$
are not smooth at the origin.
\item For the Bessel functions, $K_0(x)=-\ln\frac x2\,I_0(x)+(\text{series in }x^2)$, and $I_0$ is
itself a series in $x^2$. So $K_0\,\mathcal N=(\text{analytic in }r^2)+\ln r\times(\text{analytic in }r^2)$.
The analytic part contributes nothing to any order in $1/p$ (it is not exactly zero, but beyond all powers).
Only the $\ln r$ part leaves a power-law tail. $K_1$ is the same, with an extra factor $re^{i\phi}$.
\item Physics: a logarithm in $r$ is a power tail in momentum. The $\ln r$ of $K_0$ is the
perturbative $1/k^2$ tail of the $\gamma\to c\bar c$ splitting, and the $\ln\frac1{r\Lambda}$ of MV
is the $1/k^4$ one-gluon tail of the target. A smooth piece such as the GBW $r^2$ (a Gaussian in $k$)
carries no hard momentum. If neither the wavefunction nor the target supplies the large $p$, the
quark--antiquark cancellation (colour transparency) removes the contribution. The logarithms are the
only hard sources, so they alone set the power.
\end{itemize}

We need the derivatives at those points. Writing $G(a)=1/\Gamma\big(\tfrac{1+\nu-a}{2}\big)$ and using
$1/\Gamma(-n+\delta)=(-1)^n n!\,\delta+O(\delta^2)$ and, for $n=1$,
$1/\Gamma(-1+\delta)=-\delta+(1-\gamma_E)\delta^2$:
\begin{align*}
  \partial_aM_0(3)&=\frac{4}{p^4}, &
  \partial_a^2M_0(3)&=\frac{8}{p^4}\Big[-\ln p+\ln2+1-\gamma_E\Big],\\
  \partial_aM_1(2)&=-\frac{2}{p^3}, &
  \partial_aM_1(4)&=\frac{16}{p^5}.
\end{align*}

\section*{3. Small-$r$ expansions}

\textbf{Bessel functions} ($x=mr\ll1$):
\[
  K_0(x)=-\ln\frac{x}{2}-\gamma_E+O(x^2\ln x),\qquad
  K_1(x)=\frac1x+\frac x2\Big[\ln\frac x2+\gamma_E-\frac12\Big]+O(x^3\ln x).
\]
\textbf{Dipole amplitude.} The MV$^e$ initial condition of our dipole files is
\[
  \mathcal{N}_F(r)=1-\exp\Big[-\frac{r^2Q_{s0}^2}{4}\,\ln\Big(\frac{1}{r\Lambda}+e_c\,e\Big)\Big]
  \;\simeq\;\frac{r^2Q_{s0}^2}{4}\,\ell(r)+O\big(r^4Q_{s0}^4\big),
\]
with $\Lambda=0.241$~GeV. The $O(r^4)$ terms are the multiple scatterings (saturation); they give
corrections suppressed by $Q_{s0}^2/p^2$, which we drop. With $r_0\equiv1/(e_ce\Lambda)\approx0.081$~GeV$^{-1}$
($1/r_0\approx12$~GeV), the logarithm has two limits:
\begin{itemize}
\item $r\gg r_0$: $\ell=\ln(e_ce)+\ln\big(1+\frac{r_0}{r}\big)\simeq\ln(e_ce)+\frac{r_0}{r}-\dots$, so
\[
  \mathcal N\simeq\frac{Q_{s0}^2}{4}\Big[\ln(e_ce)\,r^2+r_0\,r+\dots\Big].
\]
The first term is GBW-like (smooth). The second, $r_0|\mathbf r|$, is \textbf{not} smooth in 2D, so by
Section 2 it gives a power-law tail. In momentum space,
$F(k)=\int dr\,r\,J_0(kr)\,(1-\mathcal N)\simeq Q_{s0}^2r_0/(4k^3)$ for $k\ll1/r_0$.
\item $r\ll r_0$: $\ell\simeq\ln\frac{1}{r\Lambda}$. This is the MV dipole, whose logarithm
is the perturbative one-gluon tail $\varphi(k)\propto1/k^4$ of the target.
\end{itemize}
So $1/r_0\approx12$~GeV is where the tail changes from $k^{-3}$ to $k^{-4}$, not where it starts. Local
power $d\ln F/d\ln k$ of the initial condition (proton, $Q_{s0}^2=0.06$~GeV$^2$, direct integration):
\begin{center}
\begin{tabular}{lccccccc}
\toprule
$k$ [GeV] & 1 & 2 & 3 & 5 & 8 & 12 & 20 \\
\midrule
MV$^e$, $e_c=18.9$ & $-9.5$ & $-5.6$ & $-3.2$ & $-3.1$ & $-3.1$ & $-3.2$ & $-3.2$ \\
MV, $e_c=1$        & $-3.5$ & $-3.6$ & $-3.6$ & $-3.7$ & $-3.8$ & $-3.9$ & $-3.9$ \\
\bottomrule
\end{tabular}
\end{center}
(GBW is a Gaussian in $k$, with no power-law tail.) The formula $Q_{s0}^2r_0/(4k^3)$ agrees with the
numerics within 5\% for $3\le k\le8$~GeV. The power-law tail of MV$^e$ starts at $k\approx3$~GeV, as in
the momentum-space plots of Lappi--M\"antysaari (arXiv:1309.6963).

\section*{4. GBW: $\mathcal{N}=Q_{s0}^2r^2/4$}

\textbf{$\nu=0$.} $r\,K_0\,\mathcal N=\frac{Q_{s0}^2}{4}\,r^3\big[-\ln r-\ln\frac m2-\gamma_E\big]$.
The $r^3$ term is smooth and gives zero; only $r^3\ln r$ survives:
\[
  I_0\simeq-\frac{Q_{s0}^2}{4}\,\partial_aM_0(3)=-\frac{Q_{s0}^2}{p^4}.
\]
\textbf{$\nu=1$.} $r\,K_1\,\mathcal N=\frac{Q_{s0}^2}{4}\big[\frac{r^2}{m}+\frac{m\,r^4}{2}\ln r+\dots\big]$.
The leading term $r^2/m$ is smooth and gives \textbf{zero}; the first non-zero term comes from the
logarithm of the \emph{next} order of $K_1$, which carries $m^2$:
\[
  I_1\simeq\frac{Q_{s0}^2\,m}{8}\,\partial_aM_1(4)=\frac{2\,m\,Q_{s0}^2}{p^5}.
\]
Both agree with the exact results for this dipole,
$I_0=-Q_{s0}^2(p^2-m^2)/(p^2+m^2)^3$ and $I_1=2mpQ_{s0}^2/(p^2+m^2)^3$.

\textbf{Cross section.}
\[
  m^2\Big[(z^2+\bar z^2)I_1^2+I_0^2\Big]\simeq m^2Q_{s0}^4\Big[\frac{1}{p^8}+\frac{4(z^2+\bar z^2)m^2}{p^{10}}\Big]
  \;\propto\;\frac{m^2Q_{s0}^4}{p^8}.
\]
So for an $r^2$ dipole the power is $-8$: the amplitude is $Q_{s0}^2/p^4$ (instead of the wavefunction
$1/p^2$) because only the $\ln(mr)$ of $K_0$ survives the cancellation, and the $I_1$ term is a further
$m^2/p^2$ smaller.

\section*{5. MV: $\mathcal{N}=\frac{Q_{s0}^2r^2}{4}\big[\ln\frac{1}{r\Lambda}+c\big]$}

Now there is a second logarithm, from the target.

\textbf{$\nu=0$.} $r\,K_0\,\mathcal N=\frac{Q_{s0}^2}{4}\,r^3\big[-\ln r-\ln\Lambda+c\big]
\big[-\ln r-\ln\frac m2-\gamma_E\big]$. The $r^3\ln^2r$ term needs $\partial_a^2M_0(3)$:
\[
  I_0\simeq\frac{Q_{s0}^2}{4}\Big[\partial_a^2M_0(3)+\Big(\ln\frac m2+\gamma_E+\ln\Lambda-c\Big)\partial_aM_0(3)\Big]
  =-\frac{Q_{s0}^2}{p^4}\Big[\ln\frac{p^2}{2m\Lambda}-2+\gamma_E+c\Big].
\]
Same power as GBW, with a logarithm: $\ln\frac{p^2}{m\Lambda}=\ln\frac pm+\ln\frac p\Lambda$, one from
$K_0$ and one from the target.

\textbf{$\nu=1$.} The leading term is now $\frac{Q_{s0}^2}{4m}\,r^2\big[-\ln r-\ln\Lambda+c\big]$. The
$r^2$ part still gives zero, but $r^2\ln r$ does not:
\[
  I_1\simeq-\frac{Q_{s0}^2}{4m}\,\partial_aM_1(2)=\frac{Q_{s0}^2}{2\,m\,p^3}.
\]
This is two powers of $p$ larger than in GBW. Physically: the target can now give a single hard kick
$k\sim p$ (the $1/k^4$ tail), and then the cancellation between quark and antiquark does not
happen.

\textbf{Cross section.}
\[
  m^2\Big[(z^2+\bar z^2)I_1^2+I_0^2\Big]\simeq
  \frac{(z^2+\bar z^2)\,Q_{s0}^4}{4\,p^6}
  +\frac{m^2Q_{s0}^4}{p^8}\Big[\ln\frac{p^2}{2m\Lambda}-2+\gamma_E+c\Big]^2 .
\]
The mass drops out of the first term. At very large $p$ it wins and the power is $-6$. The two
terms are equal at $p\simeq2m\,L/\sqrt{z^2+\bar z^2}$, with $L$ the bracket: for $z=1/2$ and $m=1.5$~GeV,
$p\approx20$--$25$~GeV.

\textbf{Numerical check} (direct integration with $\mathcal N=\frac{r^2}{4}\ln\frac1{r\Lambda}$, $m=1.5$~GeV):
\begin{center}
\begin{tabular}{lcccc}
\toprule
$p$ [GeV] & 20 & 50 & 100 & 200 \\
\midrule
$I_0$ numerical / asymptotic & 0.978 & 0.996 & 0.999 & 1.000 \\
$I_1$ numerical / asymptotic & 1.108 & 1.024 & 1.007 & 1.002 \\
\bottomrule
\end{tabular}
\end{center}

\section*{5b. BK evolution generates the logarithm}

BK evolution changes the power because it makes a dipole that is smooth at $r=0$ non-smooth there.

\textbf{One step of evolution.} Start from the GBW-like $\mathcal N=Q^2r^2/4$ (smooth, $p^{-8}$). At small
$r$ the BK equation is linear,
\[
  \partial_Y\mathcal N(r)=\frac{\bar\alpha}{2\pi}\int d^2\mathbf z\,\frac{r^2}{X^2Z^2}
  \big[\mathcal N(X)+\mathcal N(Z)-\mathcal N(r)\big],\qquad X=|\mathbf x-\mathbf z|,\ Z=|\mathbf z-\mathbf y|.
\]
The important region is $r\ll z\ll1/Q_s$: the gluon is emitted far from the small dipole. There
$X\simeq Z\simeq z$ and
\[
  \frac{r^2}{z^4}\cdot2\cdot\frac{Q^2z^2}{4}=\frac{Q^2r^2}{2z^2},\qquad
  \int_r^{1/Q_s}\frac{d^2z}{z^2}=2\pi\ln\frac1{rQ_s},
\]
so
\[
  \partial_Y\mathcal N\simeq2\bar\alpha\,\frac{Q^2r^2}{4}\ln\frac1{rQ_s}+(\text{smooth}),\qquad
  \mathcal N(Y)\simeq\frac{Q^2r^2}{4}\Big[c+2\bar\alpha Y\ln\frac1{rQ_s}\Big].
\]
This is the MV form of Section 5, with $2\bar\alpha Y$ instead of $1$ in front of the logarithm. Even
from GBW, or from MV$^e$ (whose $\ln\frac1r$ appears only at $r\ll r_0$), evolution produces the $r^2\ln r$ term, which feeds
$I_1$ without the $m^2$ suppression:
\[
  m\,I_1\to2\bar\alpha Y\,\frac{Q^2}{2p^3}\qquad\Longrightarrow\qquad
  (2\bar\alpha Y)^2\,\frac{(z^2+\bar z^2)\,Q^4}{4\,p^6}.
\]

\textbf{Physics.} A smooth $r^2$ dipole is a target with a Gaussian gluon distribution in $k$, i.e.\ no
hard gluons. Evolution emits gluons with the perturbative $1/k^2$ spectrum, which builds the $1/k^4$ tail
of the unintegrated gluon distribution ($\ln\frac1r$ in coordinate space; with repeated emissions, the
double-log growth $\mathcal N\sim r^2\,xG(x,1/r^2)$). With that tail the target can give one hard kick
$k\sim p$, and the quark--antiquark cancellation no longer happens.

\textbf{Where $r^{2\gamma}$ applies.} The geometric-scaling form $\mathcal N\propto(r^2Q_s^2)^{\gamma}$ with
$\gamma<1$ holds only for $r$ near $1/Q_s$ (extended-scaling window). At $p\sim10$~GeV with
$Q_s\lesssim1$~GeV we are far outside it, and the right description is the one above: $\gamma\simeq1$ with
generated logarithms, giving $p^{-6}\times\log^2$. A tail flatter than $p^{-6}$ would need $p$ close to $Q_s$.

\textbf{Size in our dipole (numerical).} The one-step formula suggests a small coefficient
($2\bar\alpha Y\approx0.3$ for $Y\approx1$), but it underestimates how fast rcBK builds the logarithm.
Fitting the dipole files (\texttt{data/Pb/mve/glauber\_mve\_0}, centre of Pb) to
$4\mathcal N/r^2=A+B\ln\frac1r$ for $r=0.02$--$0.2$~GeV$^{-1}$ ($p\sim5$--$50$~GeV):
\begin{center}
\begin{tabular}{lcccc}
\toprule
$x_{\mathbb P}$ & $10^{-2}$ (IC) & $4\times10^{-3}$ & $10^{-3}$ & $10^{-4}$ \\
$Y=\ln(0.01/x_{\mathbb P})$ & 0 & 1 & 2.3 & 4.6 \\
\midrule
$B$ [GeV$^2$] & 0.12 & 0.27 & 0.63 & 1.93 \\
\bottomrule
\end{tabular}
\end{center}
Pure MV with the same $Q_{s0}^2$ has $B=Q_{s0}^2=0.22$. So at $x_{\mathbb P}=10^{-3}$ the generated logarithm
is already three times the MV one. The $K_1$ term grows as $B^2$, without $m^2$, and dominates: at
\textbf{fixed} $x_{\mathbb P}\le10^{-3}$ the partonic spectrum goes as $p^{-6}$ to $p^{-6.8}$ from
$p\approx5$~GeV on, and the $K_1$ term is $40$--$70\%$ of the cross section at $8$--$12$~GeV.
In the real kinematics, however, $x_{\mathbb P}$ grows with $p$ and this
is undone (Section 7c).

\section*{6. What this means for our dipole}

The MV$^e$ logarithm $\ln(1/(r\Lambda)+e_ce)$ becomes the MV one only for $r\ll r_0$, that is
\[
  p\gg 1/r_0=e_c\,e\,\Lambda\approx18.9\times2.72\times0.241~\text{GeV}\approx12~\text{GeV}.
\]
Below that the dipole is $r^2\ln(e_ce)$ plus the non-smooth $r_0|\mathbf r|$ term of Section 3. The
latter feeds $I_1$ without the $m^2$ suppression (Eq.~\eqref{eq:mellin} with $a=1$, no derivative: $M_1(1)=1/p^2$):
\[
  m\,I_1\simeq m\,I_1^{\rm GBW}+\frac{Q_{s0}^2\,r_0}{4\,p^2},
\]
which agrees with the numerics within 3--16\% for $3\le p\le8$~GeV (27\% at 12~GeV, where $r_0p\to1$).
Compared with the $K_0$ term $m\,I_0\simeq mQ_{s0}^2\ln(e_ce)/p^4$ it is small because of $r_0$:
\[
  \frac{(z^2+\bar z^2)\,m^2I_1^2}{m^2I_0^2}\simeq\frac{(z^2+\bar z^2)\,r_0^2\,p^4}{16\,m^2\ln^2(e_ce)}
  \approx\frac{p^4}{1.7\times10^5~\text{GeV}^4}\quad(z=1/2),
\]
about 2\% at 8~GeV and 12\% at 12~GeV. So the tail is there at all $p\gtrsim3$~GeV, but the GBW-like
$m^2/p^8$ term still dominates the cross section up to about 12~GeV. Numerically, the $I_1$ fraction
$(z^2+\bar z^2)I_1^2/[(z^2+\bar z^2)I_1^2+I_0^2]$ at $z=1/2$ is $0.14$ at 8 and 12~GeV for MV$^e$, against
$0.07$ and $0.03$ for GBW. Local power
$d\ln[(z^2+\bar z^2)I_1^2+I_0^2]/d\ln p$ for $z=1/2$ (direct integration, with saturation, $Q_{s0}^2=0.216$~GeV$^2$,
about the value at the center of Pb):
\begin{center}
\begin{tabular}{lccccccc}
\toprule
$p$ range [GeV] & 3--5 & 5--8 & 8--12 & 12--20 & 20--40 & 40--80 & 80--160 \\
\midrule
GBW ($r^2$, exact)         & $-6.5$ & $-7.4$ & $-7.7$ & $-7.9$ & $-8.0$ & $-8.0$ & $-8.0$ \\
MV$^e$, $e_c=18.9$         & $-6.5$ & $-7.3$ & $-7.5$ & $-7.3$ & $-6.7$ & $-6.1$ & $-5.9$ \\
MV, $e_c=1$                & $-5.5$ & $-6.2$ & $-6.5$ & $-6.6$ & $-6.6$ & $-6.3$ & $-6.1$ \\
\bottomrule
\end{tabular}
\end{center}
\begin{itemize}
\item Below $3$--$5$~GeV no form is asymptotic: $p$ is not much larger than $m$ (for GBW the exact local
power is $4p^4/(p^4+m^4)-12p^2/(p^2+m^2)$ at $z=1/2$, which reaches $-8$ only slowly).
\item MV$^e$ is close to GBW up to about $12$~GeV (slightly flatter, from the $r_0|\mathbf r|$ term) and then
turns towards $-6$, as predicted.
\item BK evolution generates $r^2\ln\frac1r$ terms (Section 5b; an effective $\gamma<1$ only near
$Q_s$, see also Section 6 of \texttt{exclusive\_pt8.pdf}). Like the MV logarithm, this is a non-smooth
term, so it also feeds $I_1$ without the $m^2$ suppression: at fixed $x_{\mathbb P}\le10^{-3}$ the spectrum
is close to $p^{-6}$.
\item In the full $D^0$ calculation the photon flux, the $x_{\mathbb P}$ dependence and the fragmentation make
the spectrum steeper again (local power $-7.5$ to $-9.7$ at $y=0$, see \texttt{exclusive\_pt8.pdf}).
\end{itemize}

\textbf{Why the MV$^e$ $K_1$ term reaches $-6$ only at large $p$}. Three regimes:
\begin{itemize}
\item $p\lesssim10$~GeV: $m\,I_1\simeq2m^2Q_{s0}^2\ln(e_ce)/p^5+Q_{s0}^2r_0/(4p^2)$. The two are equal at
$p\approx(8m^2\ln(e_ce)/r_0)^{1/3}\approx10$~GeV; below, the mass-suppressed $r^2$ part wins (local power
$-8$ to $-10$, minimum $-8.5$ at 5~GeV).
\item $p\sim12$--$20$~GeV: $r_0p\sim1$, crossover; the power passes $-6$ at about 20~GeV as a transient.
\item $p\gg12$~GeV: for $r\ll r_0$, $\ln(\frac1{r\Lambda}+e_ce)\simeq\ln\frac1{r\Lambda}+\frac r{r_0}$. The
$r/r_0$ piece gives an $r^3$ term in $r\,K_1\mathcal N$, which is not a zero of $M_1$, so the MV result
$m\,I_1\to Q_{s0}^2/(2p^3)$ has corrections $\approx1.3/(r_0p)$. These fall slowly, and the power approaches
$-6$ from above:
\end{itemize}
\begin{center}
\begin{tabular}{lcccccccc}
\toprule
$p$ [GeV] & 5 & 10 & 20 & 40 & 80 & 160 & 320 & 640 \\
\midrule
$r_0p$ & 0.4 & 0.8 & 1.6 & 3.2 & 6.5 & 13 & 26 & 52 \\
$m\,I_1/[Q_{s0}^2/(2p^3)]$ & 1.38 & 0.66 & 0.56 & 0.67 & 0.80 & 0.89 & 0.94 & 0.97 \\
local power of $K_1$ term & $-8.2$ & $-7.3$ & $-6.0$ & $-5.5$ & $-5.6$ & $-5.8$ & $-5.9$ & $-5.9$ \\
\bottomrule
\end{tabular}
\end{center}
The scale is set by $e_c$: with $e_c=1$, $1/r_0=e\Lambda\approx0.65$~GeV and $-6$ appears at a few GeV.

\section*{7. Why we see $p^{-8}$ and not $p^{-6}$}

Four effects. (a) and (b) are about the initial condition; (c) is why the evolved dipole, which at
fixed $x_{\mathbb P}$ gives $p^{-6}$, still gives a steep spectrum in the real kinematics.

\textbf{(a) In pure MV, $p^{-6}$ is only the final asymptotic power.} Section 5 gives two terms,
\[
  m^2\Big[(z^2+\bar z^2)I_1^2+I_0^2\Big]\simeq
  \underbrace{\frac{(z^2+\bar z^2)\,Q_{s0}^4}{4\,p^6}}_{I_1}
  +\underbrace{\frac{m^2Q_{s0}^4\,L^2}{p^8}}_{I_0},
  \qquad L=\ln\frac{p^2}{2m\Lambda}-2+\gamma_E\approx5\text{--}6 .
\]
The second term carries $m^2L^2\approx60$--$80$~GeV$^2$, which is not small compared with $p^2$ until
$p$ is large. The two terms are equal at $p\approx2\sqrt2\,mL\approx20$--$25$~GeV ($z=1/2$). Below that
the $p^{-8}$ term dominates, and its $\ln^2p$ makes it look flatter: pure MV has an effective power of
about $-6.5$ for $8<p<40$~GeV and approaches $-6$ only slowly.

\textbf{(b) Our dipole is MV$^e$, and its tail is weak below 12~GeV.} For $p\ll1/r_0\approx12$~GeV the
dipoles have $r\gg r_0$, and $\mathcal N\simeq\frac{Q_{s0}^2}{4}[\ln(e_ce)\,r^2+r_0\,r]$ (Section 3). The
$r^2$ part is GBW-like and gives $m^2/p^8$ (Section 4). The $r_0|\mathbf r|$ part is a genuine tail ($k^{-3}$
in momentum space, from $k\approx3$~GeV) and feeds $I_1$ without $m^2$, but with a weight
$\propto(r_0p)^2$: it is 2--12\% of the cross section for $8\le p\le12$~GeV (Section 6). So the power stays
close to $-8$. The table of Section 6 shows this:
\begin{center}
\begin{tabular}{lccccc}
\toprule
$p$ range [GeV] & 5--8 & 8--12 & 12--20 & 20--40 & 40--80 \\
\midrule
GBW ($r^2$)          & $-7.4$ & $-7.7$ & $-7.9$ & $-8.0$ & $-8.0$ \\
MV$^e$ (our dipole)  & $-7.3$ & $-7.5$ & $-7.3$ & $-6.7$ & $-6.1$ \\
pure MV              & $-6.2$ & $-6.5$ & $-6.6$ & $-6.6$ & $-6.3$ \\
\bottomrule
\end{tabular}
\end{center}
Up to about $12$~GeV our dipole is close to GBW; the turn towards $-6$ starts above, when $r_0p\sim1$.

\textbf{(c) $x_{\mathbb P}$ grows with $p$, and the evolution goes away.} At fixed $x_{\mathbb P}\le10^{-3}$ the
evolved dipole gives about $p^{-6}$ (Section 5b). But in the exclusive process
$x_{\mathbb P}=m_Te^{-y}/(\sqrt{s_{NN}}(1-z))$ (manuscript Eq.~(9)), so each doubling of $p$ doubles
$x_{\mathbb P}$ and lowers $Y$ by $\ln2$. The coefficient $B$ of the generated logarithm drops fast with $Y$
(table of Section 5b), so the $p^{-6}$ term loses weight as $p$ grows; above $x_{\mathbb P}=0.01$
(\texttt{XPO\_FROZEN}) the dipole is the IC. Local power of the full partonic cross section with
$x_{\mathbb P}(p)$ (centre of Pb, $z=1/2$, $\sqrt{s_{NN}}=5.36$~TeV):
\begin{center}
\begin{tabular}{lccccc}
\toprule
$p$ [GeV] & 3 & 5 & 8 & 12 & 20 \\
\midrule
$y=0$: $x_{\mathbb P}$ & $1.3\times10^{-3}$ & $1.9\times10^{-3}$ & $3.0\times10^{-3}$ & $4.5\times10^{-3}$ & $7.5\times10^{-3}$ \\
$y=0$: power & $-6.0$ & $-6.7$ & $-7.7$ & $-8.2$ & $-8.4$ \\
$y=+2$: power & $-5.4$ & $-6.0$ & $-6.9$ & $-7.4$ & $-7.6$ \\
$y=-2$: power & $-6.5$ & $-6.9$ & $-7.4$ & $-7.4$ & $-7.3$ \\
\bottomrule
\end{tabular}
\end{center}
At $y=0$ the result is steeper than $-8$ above 10~GeV, although every fixed-$x_{\mathbb P}$ curve is close
to $-6$: the power is the fixed-$x_{\mathbb P}$ slope plus the $x_{\mathbb P}$ dependence. At $y=-2$,
$x_{\mathbb P}>0.01$ already above about 4~GeV and the result is the IC one ($-7.3$, as in (b)). At
$y=+2$ ($x_{\mathbb P}\sim10^{-4}$) the logarithm survives better and the $K_1$ term is still about half of
the cross section at 12--20~GeV. So this steepening is already partonic, not only a $D^0$-level effect.

\textbf{(d) At the $D^0$ level the spectrum is steeper still.} The local power $-8$ to $-9.7$ of the $D^0$
results at $y=0$ includes, on top of (c), the photon flux and the fragmentation.

\textbf{In short:} $p^{-8}$ is the power for a dipole $\propto r^2$. MV$^e$ has a hard tail at all
$p\gtrsim3$~GeV, but below $p\approx12$~GeV its weight is $\propto(r_0p)^2$ and the $r^2$ part dominates.
rcBK evolution builds a strong $r^2\ln\frac1r$ term and gives $p^{-6}$ at fixed small $x_{\mathbb P}$, but
since $x_{\mathbb P}\propto m_T$ the evolution fades as $p$ grows, and at $y=0$ the spectrum is steep again.

\section*{8. Comparison with Appendix A of arXiv:2207.06268}

Iancu, Mueller, Triantafyllopoulos and Wei compute exclusive dijets in DIS with massless quarks.
Their reduced amplitude for a transverse photon is (their Eq.~(A.8))
\[
  \mathcal A^l_{el}(P,\bar Q)=-\frac{i\,\bar Q\,P^l}{P_\perp}\int_0^\infty dr\,r\,J_1(P_\perp r)\,K_1(\bar Qr)\,T(r),
  \qquad \bar Q^2=\vartheta(1-\vartheta)\,Q^2 ,
\]
with $T$ the fundamental dipole. This is our $m\,I_1$ with $m\to\bar Q$: a massless quark in a virtual
photon has the wavefunction $\bar Q K_1(\bar Qr)$, a massive quark in a real photon has $mK_1(mr)$.
(For massive quarks in DIS the scale is $\epsilon^2=\vartheta(1-\vartheta)Q^2+m^2$, which is $m^2$ at $Q^2=0$.)

\textbf{The two powers of $\bar Q$ in their Eq.~(A.9).} For $P_\perp\gg Q_s$ they find
\[
  \mathcal A^l_{el}\simeq-2i\,\frac{\bar Q^2P^l}{(P_\perp^2+\bar Q^2)^3}\,Q_A^2\ln\frac{P_\perp^2}{\Lambda^2}.
\]
\begin{enumerate}
\item One $\bar Q$ is the normalization of the wavefunction, $\bar Q K_1$. For us it is the $m$ in
$m\,I_1$.
\item The other comes from the cancellation (Section 2): the leading term $1/(\bar Qr)$ of $K_1$ gives
a smooth $r^2$ term that integrates to zero, and the first surviving term,
$\tfrac{\bar Qr}{2}\ln(\bar Qr)$, carries one more $\bar Q$. For us this is $I_1\simeq2mQ_{s0}^2/p^5\propto m$.
\end{enumerate}
So their result is our Section 4 result $m\cdot2mpQ_{s0}^2/(p^2+m^2)^3$ with $m\to\bar Q$ and
$Q_{s0}^2\to Q_A^2\ln(P_\perp^2/\Lambda^2)$ (they keep only the log-enhanced piece).

\textbf{Why $P^{-6}$ in DIS and $p^{-8}$ for us.}
\begin{itemize}
\item In DIS, $\bar Q\sim P_\perp$ (the cross section peaks there), so
$|\mathcal A|^2\propto\bar Q^4P_\perp^2/(P_\perp^2+\bar Q^2)^6\sim P_\perp^{-6}$. The factors $\bar Q^2$ set the
scale and are not a suppression: this is their higher-twist $P^{-6}$.
\item For a real photon, $\bar Q\to m\ll p$, and the same term is $m^2\cdot4m^2Q_{s0}^4/p^{10}$: an extra
$m^2/p^2$ on top of $p^{-8}$.
\item Massive quarks have in addition the same-helicity term $mK_0(mr)$ (our $I_0$), absent for massless
quarks and therefore absent in their appendix. It gives $m^2Q_{s0}^4/p^8$ and dominates for an $r^2$
dipole: this is our $p^{-8}$.
\end{itemize}

\textbf{Their $\bar Q\to0$ remark is our MV result.} In Eqs.~(A.12)--(A.13) they note that the
log-enhanced piece vanishes as $\bar Q^2\to0$, and that the remaining piece gives
\[
  \mathcal A^l_{el}(P_\perp,\bar Q=0)\simeq-\frac{iQ_A^2P^l}{P_\perp^4},\qquad
  |\mathcal A_{el}|^2\simeq\frac{Q_A^4}{P_\perp^6}.
\]
This is our Section 5 result $m\,I_1\to Q_{s0}^2/(2p^3)$, i.e.\ $(z^2+\bar z^2)Q_{s0}^4/(4p^6)$ in the cross
section. The factor 2 is a convention: their MV exponent has $\ln\frac{1}{r^2\Lambda^2}=2\ln\frac1{r\Lambda}$.
The physics is the same: a hard kick from the target (the MV logarithm) removes the $\bar Q$ or $m$
suppression.

\textbf{Conclusion.} The two calculations agree. Their $\bar Q^2$ factors are our $m$ factors of the $K_1$
term. Our $p^{-8}$ comes from the mass-induced $K_0$ term, which a massless calculation does not have.
We do not see the $p^{-6}$ of their Eq.~(A.13) because below about $12$~GeV the tail of MV$^e$ is the
weak $r_0|\mathbf r|$ term, with weight $\propto(r_0p)^2$ (Sections 3, 6 and 7).

\section*{9. Remark on the manuscript footnote}

The statement ``the higher-twist $p^{-6}$ of exclusive dijets in DIS is further suppressed by
$m^2/p^2$ for a real photon'' is correct for the part of the dipole that goes as $r^2$ (Section 4).
For a dipole with the perturbative tail (MV logarithm, or the logarithm generated by BK evolution) the $I_1$ term
gives $Q_{s0}^4/p^6$ without $m^2$ at very large $p$ (Section 5). Below $1/r_0$ the MV$^e$ tail also
gives an unsuppressed $I_1$, but with weight $(r_0p)^2$ (Section 6). For the IC this is a 10\%
effect in our range, but with rcBK evolution the unsuppressed $K_1$ term is 25--50\% of the partonic cross
section at $y=0$ for $5\le p\le20$~GeV, and more at forward $y$ (Section 7c). So the footnote should say ``parametrically'' or ``for $\pt$ below
$e_ce\Lambda$''.

\section*{10. Which parameters would give $p^{-6}$}

The power is fixed by the shape of $\mathcal N(r)$ at $r\sim1/p$ (Section 2), not by the run settings
of \texttt{run\_scripts/config.sh} (channel, flux, fragmentation, VEGAS calls). The shape comes from the
dipole files, which are \texttt{rbk} solutions (see the header of \texttt{data/Pb/mve/glauber\_mve\_0})
of the initial condition
\[
  \mathcal N(r)=1-\exp\Big[-\frac{(r^2Q_{s0}^2)^{\gamma}}{4}\ln\Big(\frac{1}{r\Lambda}+e_c\,e\Big)\Big].
\]

\textbf{(a) Initial condition: this is what decides it.}
\begin{center}
\begin{tabular}{lp{0.3\textwidth}p{0.45\textwidth}}
\toprule
parameter & our value & effect on the tail \\
\midrule
$e_c$ & 18.9 (\texttt{mve}); 6.5--60 in the BK posterior & sets $r_0=1/(e_ce\Lambda)$: weak $r_0|\mathbf r|$
tail (weight $\propto(r_0p)^2$) below $1/r_0$, MV logarithm above ($1/r_0\approx12$~GeV for 18.9, 4--39~GeV
for the posterior) \\
$\gamma$ & 1 (all posterior samples) & $\gamma<1$: $r^{2\gamma}$ is non-smooth,
$m^2I_1^2\propto p^{-(2+4\gamma)}$ \\
$Q_{s0}^2$ & 0.051--0.085 GeV$^2$ (posterior) & normalization only, no change of power \\
\bottomrule
\end{tabular}
\end{center}
With $e_c=1$ (pure MV) the logarithm is active above $e\Lambda\approx0.7$~GeV and the partonic power is
about $-6.5$ in our range (table of Section 6). With $\gamma<1$ the $I_1$ term falls as $p^{-(2+4\gamma)}$,
i.e.\ $-6$ or flatter. Both need new dipole files: run \texttt{rbk}
(\url{https://github.com/hejajama/rcbkdipole}) with the same options as in the header, and then the
Glauber step for the nucleus. To test the power alone, the proton file \texttt{data/proton/mve.dat} is
enough, since the power does not depend on $b$.

\textbf{(b) BK evolution, through $x_{\mathbb P}$.} Above \texttt{XPO\_FROZEN}$=0.01$
(\texttt{src/def.hpp}) the dipole is the initial condition itself. Below it, the rcBK evolution generates the
$r^2\ln\frac1r$ tail (Section 5b), with coefficient $B(Y)$ (table of Section 5b). Smaller $x_{\mathbb P}$
($x_{\mathbb P}\propto e^{-y}$, so larger $y$, or
\texttt{charm\_fixed\_qp\_xpom} at small fixed $x_{\mathbb P}$) therefore pushes towards $p^{-6}$. With
fixed $x_{\mathbb P}\le10^{-3}$ the effect is large ($p^{-6}$ to $p^{-6.8}$, Section 5b), but with
$x_{\mathbb P}(p)$ from the kinematics it fades as $p$ grows (Section 7c): forward rapidity, where
$x_{\mathbb P}$ stays small, is where it shows.

\textbf{(c) The $p$ range.} \texttt{PT\_VALS} stops at 12~GeV. For the IC the turn towards $-6$ starts only
above that ($-6.7$ for 20--40~GeV, $-6.1$ for 40--80~GeV). This holds at fixed $x_{\mathbb P}\ge0.01$; with
$x_{\mathbb P}(p)$ at $y=0$ the large-$p$ points sit at $x_{\mathbb P}>0.01$ (IC) and the evolved
small-$p$ points are flatter, so the measured slope is steeper (Section 7c).

\textbf{Suggested test.} Run \texttt{charm\_fixed\_qp\_xpom} at a fixed $x_{\mathbb P}$ ($10^{-2}$ and
$10^{-3}$) up to $p\approx80$~GeV and plot $p^6\,d\sigma$: a $p^{-6}$ tail is a plateau.

\textbf{Caveat.} $e_c$ and $\gamma$ come from the fit to HERA data. Changing them is a diagnostic of where
the $p^{-6}$ comes from, not another prediction, unless there is a physical reason (e.g.\ a refit with
$\gamma$ free). And at the $D^0$ level the flux, $x_{\mathbb P}$ and fragmentation steepen the spectrum
(Section 7c), so a partonic $p^{-6}$ will not look like $p^{-6}$ in the $D^0$ figure.

\section*{11. $R_{pA}$ and the manuscript's $\mathcal N\propto(r^2Q_s^2)^\gamma$}

In the $R_{pA}$ discussion (Sec.~IV D, after Eq.~(38)) the manuscript says that BK evolution reduces
$N_A/N_p$ as $x_{\mathbb P}$ decreases, because ``for small dipoles $N\propto(r^2Q_s^2)^\gamma$ with
$\gamma<1$, so the ratio becomes $(Q_{s,A}^2/Q_{s,p}^2)^\gamma$''. This is about the $A$ dependence, not the
$p$ dependence of Sections 2--7. (In the initial condition, Eqs.~(ICBK), $\gamma=1$ is fixed: that
$\gamma$ only labels the MV/MV$^\gamma$/MV$^e$ family of Lappi--M\"antysaari.)

\textbf{At $x_0$.} In the dilute limit $N_A=T_AS_\perp N_p$, so $N_A/N_p\propto A^{1/3}$ independent of $r$,
and $R_{pA}\approx2.5$ (Eq.~(38)).

\textbf{Evolution lowers the ratio, through saturation.} From Section 5b,
$\partial_YN\simeq2\bar\alpha N\ln\frac1{rQ_s}$. The logarithm comes from daughter dipoles with
$r<z<1/Q_s$: above $1/Q_s$ they are saturated and stop contributing. The nucleus has the larger $Q_s$, so
its logarithm is shorter and its small dipoles grow more slowly:
\[
  \partial_Y\ln\frac{N_A}{N_p}\simeq-2\bar\alpha\ln\frac{Q_{s,A}}{Q_{s,p}}=-\bar\alpha\ln\frac{Q_{s,A}^2}{Q_{s,p}^2}<0.
\]
Purely linear (BFKL) evolution would keep the ratio fixed; the suppression comes from the larger
nuclear saturation scale.

\textbf{Asymptotically: geometric scaling.} After enough evolution $N\simeq f\big(r^2Q_s^2(Y)\big)$ for both
targets, and for $r$ somewhat below $1/Q_s$, $N\simeq(r^2Q_s^2)^{\gamma_s}$ with $\gamma_s\approx0.6$--$0.8$.
Then
\[
  \frac{N_A}{N_p}=\Big(\frac{Q_{s,A}^2}{Q_{s,p}^2}\Big)^{\gamma_s}<\frac{Q_{s,A}^2}{Q_{s,p}^2}\sim A^{1/3}:
\]
the nucleus keeps only part of its naive $A^{1/3}$ enhancement (the same mechanism as the forward
suppression of $R_{dAu}$).

\textbf{Consistency with Section 5b.} The $r^{2\gamma}$ form holds only near $1/Q_s$. The $R_{pA}$ argument is
made at small $p$, where $r\sim1/m_c$ is not far from $1/Q_s$, so it applies there. For $r\ll1/Q_s$ (large
$p$), $\gamma\to1$ with logarithms and the ratio returns towards $A^{1/3}$, which agrees with the growth of
$R_{pA}$ with $p$ described in the manuscript.

\textbf{Caveats for the wording.}
\begin{itemize}
\item ``For small dipoles'' is too broad: it is for dipoles not much smaller than $1/Q_s$.
\item $\gamma<1$ is the asymptotic form. With $Y=\ln(0.01/x_{\mathbb P})\sim1$--$3$ and running coupling the
solution has not fully reached scaling; the derivative formula above describes the early evolution
better.
\end{itemize}
Possible wording: ``In addition, BK evolution reduces $N_A/N_p$ as $x_{\mathbb P}$ decreases, because the
larger nuclear saturation scale limits the growth of small dipoles; asymptotically, for $r_\perp$ not much
smaller than $1/Q_s$, $N\propto(r_\perp^2Q_s^2)^\gamma$ with $\gamma<1$ and the ratio becomes
$(Q_{s,A}^2/Q_{s,p}^2)^\gamma$.''

\textbf{Footnote 1 of the manuscript} (``This is because..'' is unfinished). From Section 8, a possible
completion: ``This is because in DIS the factors $\bar Q^2$ are of order $P_\perp^2$, while for a real photon
they are replaced by $m^2\ll p_\perp^2$. The $m^2/p_\perp^8$ behaviour comes from the term with equal quark
helicities, $mK_0(mr_\perp)$, which is absent for massless quarks.'' (Also: ``that the one'' $\to$ ``than
the one''.)

\end{document}
